% Source: https://artofproblemsolving.com/community/c3h2220396_hard_problem
% Stable identity: aops-topic-2220396-q01
Let the numbers be $x$ and $y$. The pair $(x,y)$ is uniform in the unit square, so probabilities are areas. The condition $\dfrac{x+y}{2}\le\dfrac{2}{3}$ means $x+y\le\dfrac{4}{3}$.

\begin{asy}
usepackage("amsmath");
size(4.5cm);
defaultpen(fontsize(10pt));
fill((0,0)--(1,0)--(1,1/3)--(1/3,1)--(0,1)--cycle,lightblue);
draw((0,0)--(1,0)--(1,1)--(0,1)--cycle);
draw((1/3,1)--(1,1/3),blue+linewidth(0.9));
label("$x+y=\dfrac{4}{3}$",(1,0.8),E,blue+fontsize(8pt));
label("$\dfrac{1}{3}$",(1,1/3),E,fontsize(8pt)); label("$\dfrac{1}{3}$",(1/3,1),N,fontsize(8pt));
label("$0$",(0,0),SW); label("$1$",(1,0),S); label("$1$",(0,1),W);
label("$x$",(0.5,0),S); label("$y$",(0,0.5),W);
\end{asy}

The line $x+y=\dfrac{4}{3}$ cuts off the corner triangle near $(1,1)$ with vertices $\left(\dfrac{1}{3},1\right)$, $(1,1)$, $\left(1,\dfrac{1}{3}\right)$. Its legs are $\dfrac{2}{3}$, so its area is $\dfrac{1}{2}\cdot\left(\dfrac{2}{3}\right)^2=\dfrac{2}{9}$. Everything else in the square satisfies the condition, so the probability is
\[1-\dfrac{2}{9}=\boxed{\dfrac{7}{9}}.\]
